\uuid{nyeC}
\titre{Ensembles}
\auteur{bodin ; croizat ; sacre}
\organisation{exo7}
\licence{CC BY-NC-SA 4.0}
\datecreate{2022-03-08}
\niveau{L1}
\module{Algèbre}
\chapitre{Logique, ensemble, raisonnement}
\sousChapitre{Ensemble}
\difficulte{2}
\theme{Ensembles}
\exo7qcmid{hilisit.logique.028}

\contenu{
\begin{question}
\qclassification{100.02}
Soient $f,g,h$ des fonctions définies par
$$f(x) = x^2-3x \qquad g(x) = 2x+1 \qquad h(x) = \frac{x}{x-1}.$$
Quelles sont les affirmations vraies ?
\begin{answers}
	\bad{$(f \circ g)(x) =4x^2-2x+4$}
	\bad{$(g \circ f)(x) = 2x^2+5$}
	\bad{$(h \circ f)(x) = \left(\frac{x}{x-1}\right)^2 - \frac{3x}{x-1}$}
	\good{$(g \circ h)(x) = \frac{3x-1}{x-1}$}
\end{answers}
\begin{explanations}
$(f \circ g)(x) = f(g(x)) = (2x+1)^2-3(2x+1) = 4x^2-2x-2$, \\
$(g \circ f)(x) = g(f(x)) = 2(x^2-3x)+1 = 2x^2-6x+1$. \\
$(h \circ f)(x) = h(f(x)) = \frac{x^2-3x}{x^2-3x-1}$ ce qui est proposé c'est $(f \circ h)(x)$.\\
$(g \circ h)(x) = g(h(x)) = 2\frac{x}{x-1} +1 = \frac{2x}{x-1} + \frac{x-1}{x-1} = \frac{3x-1}{x-1}$.
\end{explanations}
\end{question}
}
